Let the inradius be $r$, and let $\gamma = \dfrac{C}{2}$. Since $D$ lies on segment $AI$, $ID=r$ and the tangent at $D$ is perpendicular to $AI$. The angle between $BI$ and $AI$ is $90^\circ+\gamma$, so its intersection with that tangent gives $DX=r\cot\gamma$. The exradius identity $\dfrac{r_a}{r} = \cot (\dfrac{B}{2}) \cot (\dfrac{C}{2})$ follows from $\dfrac{r_a}{r} = \dfrac{s}{s-a}$ and the standard half-angle formulas. With $B=60^\circ$, it gives $r_a=\sqrt3\,r\cot\gamma=\sqrt3\,DX$. Since $AD=r_a$, the right triangle $ADX$ has $\tan \angle X A I = \dfrac{D X}{A D} = \dfrac{1}{\sqrt{3}}$. Therefore $\boxed{\angle XAI=30^\circ}$.
\begin{asy}
import graph;
usepackage("amsmath");
import geometry;
size(7cm);
pair I=(0,0),D=(1,0),X=(1,-1.2),A=(1+1.2sqrt(3),0);draw(circle(I,1));draw(I--A--X--I);draw((1,-1.5)--(1,1.5));label("$I$",I,(-1,0));label("$D$",D,(0,1));label("$A$",A,(1,0));label("$X$",X,(0,-1));label("$AD=r_a$",(A+D)/2,(0,1));
\end{asy}

